%%%%%%%%%%%%%%%%%%%%%%%%%%%%%%%%%%%%%%%%%%%%%%%%%%%%%%%%%%%%%%%%%%%%%%%%%%%%%%%
%% Seção 4 --- Resultados e discussão
%%
%% Mesma ordem da seção de métodos: aquisição, estimação, ensaio de bancada e
%% comparação com a literatura. A origem de cada resultado (sinal sintetizado,
%% dado real decimado ou registro de bancada) é informada no próprio texto e nas
%% notas das tabelas e figuras.
%%
%% Fontes dos valores:
%%   sinal sintetizado    experiments/figures/data/fig1*.csv e fig2*.csv
%%                        (scripts/export_source_data.py)
%%   dado real decimado   experiments/figures/data/tab_deteccao_*.csv
%%                        (scripts/export_deteccao_data.py)
%%   bancada              logs/device-monitor-261011-155502.log e
%%                        logs/c3-device-validation-20261024-183820.log
%%%%%%%%%%%%%%%%%%%%%%%%%%%%%%%%%%%%%%%%%%%%%%%%%%%%%%%%%%%%%%%%%%%%%%%%%%%%%%%

\makeatletter
%% booktabs não é carregado pelo template; se ausente, as réguas viram \hline
\@ifundefined{toprule}{%
  \newcommand{\toprule}{\hline}%
  \newcommand{\midrule}{\hline}%
  \newcommand{\bottomrule}{\hline}%
}{}
\@ifundefined{emandamento}{%
  \@ifundefined{textcolor}%
    {\newcommand{\emandamento}[1]{#1}}%
    {\newcommand{\emandamento}[1]{\textcolor{red}{#1}}}%
}{}
\makeatother

%% Caminhos para as figuras geradas pelo pipeline do repositório.
%% NÃO copiar as figuras para dentro do TCC: elas são regeneráveis por
%% experiments/figures/ e uma segunda cópia divergiria em silêncio.
%% Ajustar o prefixo abaixo se a pasta de compilação mudar de lugar.
\providecommand{\kxfig}[1]{../../experiments/figures/output/#1}

\chapter{Resultados e discussão}
\label{cap:resultados}

Em dado real decimado, a detecção foi de \SI{75,0 \pm 3,5}{\percent} para desbalanceamento
e de no máximo \SI{20,7}{\percent} para defeito de rolamento sem massa de desbalanceamento
sobreposta. Com o limiar padrão, o falso alarme foi de \SI{31,2 \pm 4,2}{\percent}, contra
\SI{2,0 \pm 1,3}{\percent} com o limiar calibrado.

\section{Origem dos resultados}
\label{sec:proc}

Os resultados vêm de sinal sintetizado, de dado real decimado e de registros de bancada.

\textbf{Sinal sintetizado.} O sinal é gerado pelo modelo da seção~\ref{sec:dev-sinal}. Ele
        permite comparar métodos, mas não indica o desempenho em uma máquina.

\textbf{Dado real decimado.} As gravações vêm de uma bancada instrumentada e foram decimadas
        em \emph{software} para a taxa do dispositivo. O sinal é medido, mas a aquisição é
        simulada.

\textbf{Registro de bancada.} \emandamento{Os valores vêm do próprio dispositivo, sem motor
        e sem falha induzida. Eles permitem avaliar a aquisição e a transmissão, mas não a
        detecção.}

\section{Limites da aquisição}
\label{sec:res-aquisicao}

Com amostragem a \SI{1}{\kilo\hertz}, a ressonância em que aparece a energia do defeito de
rolamento fica fora da banda. Em dado real decimado, na condição sem massa sobreposta, o
desbalanceamento foi a única classe com detecção acima de \SI{20,7}{\percent}.

\subsection{Banda e evidência de defeito de rolamento}

A taxa de \SI{1}{\kilo\hertz} põe a frequência de Nyquist em \SI{500}{\hertz}. A frequência
característica de defeito de pista externa, de \SI{104,6}{\hertz} na geometria considerada,
fica dentro dessa banda, mas a ressonância de \SI{3,5}{\kilo\hertz} do modelo de sinal fica
fora.

Em sinal sintetizado, o contraste entre as condições com e sem defeito foi de \num{413}
vezes por análise de envelope na banda de ressonância, amostrada a \SI{20}{\kilo\hertz}, e
de \num{20} vezes por espectro direto a \SI{1}{\kilo\hertz} (Figura~\ref{fig:banda}). As
duas medidas diferem na banda e no método de demodulação, então esse par não permite dizer
quanto da perda vem da banda.

Sem filtro abaixo de $f_s/2$, a ressonância de \SI{3,5}{\kilo\hertz} reaparece rebatida em
\SI{500}{\hertz}, onde não há nenhuma componente física. A \SI{1}{\kilo\hertz} com filtro, a
curtose quase não mudou com o defeito, de \num{1,76} para \num{1,77}. Sem filtro, ela foi de
\num{1,82} para \num{4,55}, valor próximo do obtido a \SI{20}{\kilo\hertz}. Assim, um detector
pode separar as classes por causa do rebatimento e passar na validação. O MPU6050 sempre opera
com o filtro interno, mas o fabricante não informa a atenuação desse filtro acima da banda
(seção~\ref{sec:fund-mpu6050}), e este trabalho não a mediu.

\begin{figure}[!htb]
\caption{Efeito do limite de banda sobre a evidência de defeito de rolamento}
\label{fig:banda}
\centering
\includegraphics[width=\textwidth]{\kxfig{fig1_banda_sensor.pdf}}
\fonte{Elaboração própria.}
\legenda{Sinal sintetizado. \textbf{a}, Densidade espectral de potência para rolamento sadio e com defeito de pista externa, com eixo vertical em escala logarítmica. A faixa clara vai até \SI{500}{\hertz}, e a escura, até os \SI{184}{\hertz} do filtro usado no \emph{firmware}. \textbf{b}, Razão entre as amplitudes na frequência característica nas duas condições. \textbf{c}, Razão entre as condições para quatro descritores de impacto, calculados sobre \SI{2}{\second} de sinal; a linha tracejada marca a razão 1, sem separação. \textbf{d}, Frequência aparente da ressonância após subamostragem sem filtro.}
\end{figure}

\subsection{Detecção por modo de falha}

Em dado real decimado, com o limiar calibrado para falso alarme alvo de \SI{1}{\percent}, o
falso alarme medido foi de \SI{2,0 \pm 1,3}{\percent} (média e desvio-padrão entre as cinco
partições). A Tabela~\ref{tab:deteccao} mostra a detecção por classe nesse ponto de
operação.

\begin{table}[!htb]
\IBGEtab{%
\caption{Detecção por modo de falha a \SI{2}{\percent} de falso alarme medido, em dado real decimado}%
\label{tab:deteccao}%
}{%
\begin{tabular}{lc}
\toprule
Modo de falha & Detecção (\si{\percent}) \\
\midrule
Desbalanceamento                      & \num{75,0 \pm 3,5} \\
Rolamento \emph{overhang} (puro)      & \num{20,7 \pm 7,0} \\
Rolamento \emph{underhang} (puro)     & \num{14,8 \pm 4,3} \\
Desalinhamento vertical               & \num{7,2 \pm 4,3} \\
Desalinhamento horizontal             & \num{5,0 \pm 3,1} \\
\bottomrule
\end{tabular}%
}{%
\fonte{Elaboração própria, com dados de \citeonline{mafaulda}.}%
\nota{A incerteza é o desvio-padrão entre as cinco partições. As duas classes de rolamento correspondem à condição sem massa de desbalanceamento sobreposta.}%
}
\end{table}

Nas duas classes de desalinhamento, o desvio entre partições é da ordem do próprio valor, e
\SI{7,2 \pm 4,3}{\percent} admite valores próximos de zero. Por isso, conclusões baseadas
nesses dois números são fracas (seção~\ref{sec:res-discussao}).

As classes de rolamento foram ensaiadas com \SI{0}{\gram}, \SI{6}{\gram}, \SI{20}{\gram} e \SI{35}{\gram} de massa de
desbalanceamento sobreposta, e a detecção foi maior nas massas maiores. Na classe
\emph{underhang}, ela passou de \SI{14,8}{\percent} sem massa para \SI{87,3}{\percent} com
\SI{35}{\gram}. Na classe \emph{overhang}, passou de \SI{20,7}{\percent} para
\SI{78,3}{\percent}. Os agregados das quatro condições são de \SI{55,9}{\percent} e
\SI{42,2}{\percent}, mas esses números não podem ser lidos como detecção de defeito de
rolamento, porque o modelo parece responder à massa adicionada.

\section{Decisões de estimação}
\label{sec:res-estimacao}

Quatro decisões alteraram os resultados: o método de integração, o limiar, a métrica e o
protocolo de validação (Figura~\ref{fig:metodologia}). Exceto onde indicado, os valores desta
seção vêm de sinal sintetizado.

\begin{figure}[!htb]
\caption{Efeito de quatro escolhas metodológicas sobre os resultados}
\label{fig:metodologia}
\centering
\includegraphics[width=\textwidth]{\kxfig{fig2_metodologia.pdf}}
\fonte{Elaboração própria.}
\legenda{Sinal sintetizado. \textbf{a}, Velocidade eficaz por cinco variantes de integração, comparada com o valor analítico (linha tracejada), com soma acumulada sem remoção da média nas variantes no tempo. \textbf{b}, Velocidade sob viés de \SI{0,02}{\meter\per\second\squared}. \textbf{c}, Acurácia de um Random Forest supervisionado sobre os mesmos dados, variando a forma de partição; as barras de erro são o desvio-padrão ($n = 10$ sorteios e $n = 5$ partições), e os pontos, as partições individuais. \textbf{d}, Fração das janelas normais do ajuste classificada como anomalia, por valor do parâmetro de contaminação. \textbf{e}, Curvas ROC (característica de operação do receptor) em defeito incipiente; a faixa marca falso alarme de até \SI{10}{\percent}, usado na área parcial (pAUC).}
\end{figure}

\subsection{Integração de aceleração para velocidade}

A norma define as zonas de severidade em velocidade eficaz, e o acelerômetro entrega
aceleração, então o sinal precisa ser integrado. A Tabela~\ref{tab:res-integracao} compara os
métodos contra uma senoide de \SI{3,0}{\meter\per\second\squared} a \SI{50}{\hertz}, cuja
integral é conhecida e tem valor eficaz de \SI{6,7524}{\milli\meter\per\second}.

\begin{table}[!htb]
\IBGEtab{%
\caption{Velocidade eficaz por variante de integração, em relação ao valor analítico}%
\label{tab:res-integracao}%
}{%
\begin{tabular}{lcc}
\toprule
Método & Resultado (\si{\milli\meter\per\second}) & Erro (\si{\percent}) \\
\midrule
Soma acumulada, sem filtro passa-alta               & \num{11,69} & $+\num{73}$ \\
Soma acumulada, com filtro passa-alta               & \num{17,31} & $+\num{156}$ \\
\textbf{Integração na frequência, banda normativa}  & \num{6,75}  & \textbf{\num{0}} \\
Na frequência, com viés de \SI{0,02}{\meter\per\second\squared} & \num{6,75} & \num{0} \\
Soma acumulada, sem filtro, com o mesmo viés        & \num{55,08} & $+\num{716}$ \\
\bottomrule
\end{tabular}%
}{%
\fonte{Elaboração própria.}%
\nota{Sinal sintetizado, a \SI{5}{\kilo\hertz} durante \SI{4}{\second}. Erro em relação ao valor analítico, arredondado para inteiro. A soma acumulada parte de zero, sem remoção da média.}%
}
\end{table}

A soma acumulada parte de zero e deixa a constante de integração na média da velocidade, de
\SI{9,55}{\milli\meter\per\second}. Daí vem o erro: $\sqrt{\num{6,75}^2 + \num{9,55}^2} =
\SI{11,69}{\milli\meter\per\second}$. Com a média removida, o erro cai para
\SI{0,02}{\percent} sem filtro e para \SI{1,6}{\percent} com filtro. Um viés de
\SI{0,02}{\meter\per\second\squared}, contido na especificação do sensor, levou a soma a
\SI{55,08}{\milli\meter\per\second}, e a remoção da tendência linear também trouxe esse valor
de volta a \SI{6,75}{\milli\meter\per\second}. Sem esse cuidado, o sinal seria classificado na
zona D em qualquer máquina, enquanto o valor correto fica abaixo do limite entre as zonas C e
D em três das quatro combinações da norma.

A integração na frequência deu \SI{6,7525}{\milli\meter\per\second}, a menos de
\SI{0,01}{\percent} do valor analítico, com e sem o viés, porque o termo contínuo fica fora da
banda normativa. Nesse ensaio, os \SI{50}{\hertz} coincidem com um coeficiente exato da DFT, e
os \SI{4}{\second} de registro são mais longos que a janela de \SI{0,512}{\second} do
dispositivo.

\subsection{Limiar e falso alarme}

Em dado real decimado, avaliado em arquivos fora do treino, o limiar padrão de \num{0,5} marcou
\SI{31,2 \pm 4,2}{\percent} da operação normal como anomalia, e o limiar calibrado marcou
\SI{2,0 \pm 1,3}{\percent}.

Em sinal sintetizado, na cadeia de \SI{20}{\kilo\hertz}, o valor padrão do parâmetro de
contaminação marcou \SI{42,1}{\percent} das \num{140} janelas normais, e o valor \num{0,01}
marcou \SI{1,4}{\percent}. Essas taxas foram medidas nas próprias janelas do ajuste. Quando um
ensaio normal ficou fora do treino, os dois valores marcaram \SI{100}{\percent} das janelas nas
cinco rodadas. Os ensaios sintetizados têm ganhos diferentes, e o limiar ajustado em quatro deles
não serviu para o quinto.

Nos cinco valores testados, os escores foram idênticos, e só o deslocamento do limiar mudou. A
detecção ficou em \SI{100}{\percent} e a área sob a curva ROC em \num{1,000}, porque as classes
estão perfeitamente separadas nesse sinal, o que é improvável em dado real. O valor de
\SI{42}{\percent} vale só para este sinal.

\subsection{Métrica global e métrica no ponto de operação}

Em defeito incipiente, a área sob a curva completa caiu de \num{1,000} para \num{0,963} entre as
cadeias de \SI{20}{\kilo\hertz} e de \SI{1}{\kilo\hertz}, uma queda de \SI{3,7}{\percent}. A área
parcial caiu de \num{0,999} para \num{0,837} com limite de \SI{10}{\percent} de falso alarme, uma
queda de \SI{16,2}{\percent}, e de \num{0,998} para \num{0,795} com limite de \SI{5}{\percent}, uma
queda de \SI{20,3}{\percent}. Na região de baixo falso alarme, a queda relativa foi de \num{4,4} a
\num{5,5} vezes a indicada pela área global.

A cadeia de \SI{20}{\kilo\hertz} usa nove descritores e a de
\SI{1}{\kilo\hertz} usa sete, então a queda mistura o efeito da banda e o dos descritores. As áreas
foram calculadas sobre as janelas normais do ajuste e as janelas com defeito. Pela forma do desafio
de referência, dividida só por $\beta$, as quedas seriam de \SI{30,8}{\percent} e
\SI{39,7}{\percent}, o que muda os valores, mas não o sentido da comparação.

\subsection{Partição por janela e partição por ensaio}

Um Random Forest supervisionado, sobre os mesmos dados e os mesmos descritores, atingiu
\SI{99,0 \pm 1,0}{\percent} de acurácia sob partição aleatória por janela e
\SI{63,2 \pm 19,2}{\percent} sob partição por ensaio, uma diferença de aproximadamente \num{36}
pontos percentuais. Na partição por ensaio, três das cinco partições ficaram próximas de
\SI{50}{\percent}, o nível do acerto ao acaso, e uma alcançou \SI{100}{\percent}. As janelas têm
\SI{50}{\percent} de sobreposição, o que favorece a partição aleatória, porque janelas quase iguais
ficam no treino e no teste.

\section{Ensaio de bancada}
\label{sec:res-bancada}

\emandamento{No registro de 11 de outubro de 2026, com um ESP32-S3 como transmissor e um ESP32-C3
como receptor, o receptor aceitou \num{270} pacotes em \SI{62}{\minute}, todos com CRC de 16 bits
correto. Houve um erro de CRC no próprio rádio, que corresponde ao único pacote perdido. A contagem de
sequência recomeçou em 1 uma vez, e o receptor ficou \SI{13,4}{\minute} sem pacotes nesse intervalo. A
intensidade do sinal recebido ficou entre \SI{-73}{dBm} e \SI{-55}{dBm}, e a relação sinal-ruído,
entre \SI{9,5}{\decibel} e \SI{14}{\decibel}. Antes do primeiro pacote, o receptor registrou uma
instabilidade no barramento SPI e reiniciou o rádio. A distância entre os módulos não foi registrada.}

\emandamento{No registro de 24 de outubro de 2026, com um ESP32-C3 como transmissor, os três ciclos
gravados terminaram sem erro. A taxa de amostragem estimada ficou entre \SI{999,78}{\hertz} e
\SI{999,88}{\hertz}, com desvio máximo de \SI{0,022}{\percent}, abaixo do limite de
\SI{2}{\percent}. O intervalo entre transmissões foi de \SI{11,0}{\second}, e a fase de transmissão
durou cerca de \SI{195}{\milli\second}, próximo dos \SI{185}{\milli\second} calculados na
seção~\ref{sec:fund-lora}.}

\emandamento{Os dois registros vêm de montagens diferentes, e o segundo tem só três ciclos. Eles
indicam que a aquisição e a transmissão funcionam em bancada, mas ainda não permitem avaliar H6 por
completo. A escolha final do fator de espalhamento depende da distância até o receptor, que ainda não
foi medida.}

\section{Comparação com a literatura}
\label{sec:res-discussao}

A diferença de \num{36} pontos percentuais entre as duas formas de partição tem correlato na
literatura, que mede diferenças da ordem de dezenas de pontos percentuais em uma base pública de
rolamentos~\cite{varejao2025,vieira2026}. A concordância é de direção e de ordem de grandeza, porque
o valor deste trabalho vem de sinal sintetizado e de um classificador supervisionado.

Outros autores também separam desbalanceamento e rolamento no projeto do nó. Um nó publicado usa um
acelerômetro de banda estreita para desbalanceamento e outro de dezenas de quilohertz para
rolamento~\cite{jakobsen2024}, e o trabalho mais próximo em arquitetura demonstra desbalanceamento e
deixa rolamento para depois~\cite{kolok2025}.

O falso alarme de \SI{31,2}{\percent} com o limiar padrão, em dado real, fica próximo das
implementações embarcadas publicadas, que relatam de \SI{22}{\percent} a \SI{38}{\percent} sobre
operação normal~\cite{chevtchenko2023,kolok2025}. Esses trabalhos usam conjuntos, máquinas e limiares
diferentes, e a comparação serve só como ordem de grandeza.

Os resultados de detecção não mostram que a banda seja a causa da baixa detecção de rolamento. O
detector marca afastamento da operação normal e não identifica o tipo de defeito. Os quatro descritores
não medem a linha na frequência característica, embora o espectro direto a \SI{1}{\kilo\hertz} tenha
mostrado contraste de \num{20} vezes nessa linha em sinal sintetizado. O desalinhamento, que aparece
dentro da banda, teve detecção menor (\SI{5,0}{\percent} e \SI{7,2}{\percent}) que o defeito de
rolamento (\SI{14,8}{\percent} e \SI{20,7}{\percent}). Além disso, todas as classes usam o mesmo canal
da bancada, o sensor radial do lado \emph{underhang}, mesmo no defeito \emph{overhang}.

A baixa detecção de desalinhamento pode vir da escolha dos descritores. Sua assinatura é a razão entre
as linhas de duas vezes e uma vez a rotação, que valor eficaz, curtose, fator de crista e frequência
dominante não medem. A evidência é fraca, porque os valores têm dispersão da ordem do próprio valor. A
seção~\ref{sec:conc-futuro} indica o experimento que testaria essa explicação.
